% Part: first-order-logic
% Chapter: beyond
% Section: higher-order-logic

\documentclass[../../../include/open-logic-section]{subfiles}

\begin{document}

\olfileid{fol}{byd}{hol}

\olsection{منطق الرتب العليا}

أتاح لنا الانتقال من منطق الرتبة الأولى إلى منطق الرتبة الثانية أن
نتحدث، داخل اللغة الصورية، عن مجموعات من الأشياء في !!{domain} الرتبة
الأولى. فلماذا نتوقف عند ذلك؟ ينبغي لمنطق الرتبة الثالثة، مثلًا، أن
يتيح لنا تناول مجموعات من مجموعات الأشياء، أو لعلها حتى مجموعات تحتوي
أشياء ومجموعات من الأشياء معًا. ويتيح لنا منطق الرتبة الرابعة الحديث
عن مجموعات من الأشياء التي من ذلك النوع. وكما لعلك خمنت، يمكن تكرار
هذه الفكرة بلا حد.

يُصاغ منطق الرتب العليا في الممارسة بدلالة الدوال بدلًا من العلاقات
في كثير من الأحيان. (وهذا الفرق غير جوهري إذا أخذنا المطابقات الطبيعية
في الحسبان.) فإذا أعطيت بعض «الأصناف» الأساسية $A$، و$B$، و$C$،
و~\dots (وسنسميها الآن «أنماطًا»)، أمكننا إنشاء أنماط جديدة باشتراط
الآتي:
\begin{quote}
إذا كان $\sigma$ و$\tau$ نمطين منتهيين، كان $\sigma \to \tau$ كذلك.
\end{quote}
فكر في الأنماط بوصفها «وسومًا» تركيبية تصنف الأشياء التي نريدها في
!!{domain}نا؛ ويصف $\sigma \to \tau$ الأشياء التي تكون دوال تأخذ
أشياء من النمط~$\sigma$ إلى أشياء من النمط~$\tau$. فقد نريد، مثلًا،
نمطًا~$\Omega$ لقيمتي الصدق، «صادق» و«كاذب»، ونمطًا~$\Nat$ للأعداد
الطبيعية. وعندئذ يمكن التفكير في أشياء النمط $\Nat \to \Omega$ بوصفها
علاقات أحادية، أو مجموعات جزئية من~$\Nat$؛ وأشياء النمط $\Nat \to
\Nat$ هي دوال من الأعداد الطبيعية إلى الأعداد الطبيعية؛ وأشياء النمط
$(\Nat \to \Nat) \to \Nat$ هي «دوال دالية»، أي دوال من نمط أعلى تأخذ
الدوال إلى الأعداد.

وكما في منطق الرتبة الثانية، يمكن التفكير في منطق الرتب العليا بوصفه
نوعًا من المنطق متعدد الأصناف، فيه صنف لكل نمط من الأشياء التي نريد
النظر فيها. غير أن الأوضح عادة هو تعريف تركيب منطق الأنماط العليا من
أساسه. فيمكن، مثلًا، تعريف مجموعة من الأنماط المنتهية استقرائيًا كما
يأتي:
\begin{enumerate}
\item $\Nat$ نمط منتهٍ.
\item إذا كان $\sigma$ و$\tau$ نمطين منتهيين، كان $\sigma
  \to \tau$ كذلك.
\item إذا كان $\sigma$ و$\tau$ نمطين منتهيين، كان $\sigma \times
  \tau$ كذلك.
\end{enumerate}
يدل $\Nat$، حدسيًا، على نمط الأعداد الطبيعية، ويدل $\sigma \to
\tau$ على نمط الدوال من~$\sigma$ إلى~$\tau$، ويدل $\sigma \times
\tau$ على نمط أزواج الأشياء، أحدهما من~$\sigma$ والآخر من~$\tau$.
ويمكننا عندئذ تعريف مجموعة من الحدود استقرائيًا كما يأتي:
\begin{enumerate}
\item لكل نمط~$\sigma$ مخزون من المتغيرات $x$، و$y$، و$z$، و\dots
  من النمط~$\sigma$.
\item $\Obj 0$ حد من النمط~$\Nat$.
\item $\Obj S$ (الخلف) حد من النمط $\Nat \to \Nat$.
\item إذا كان $s$ حدًا من النمط~$\sigma$، وكان $t$ حدًا من النمط
  $\Nat \to (\sigma \to \sigma)$، كان $\Obj{R}_{st}$ حدًا من النمط
  $\Nat \to \sigma$.
\item إذا كان $s$ حدًا من النمط $\tau \to \sigma$ وكان $t$ حدًا من
  النمط~$\tau$، كان $s(t)$ حدًا من النمط~$\sigma$.
\item إذا كان $s$ حدًا من النمط~$\sigma$ وكان $x$ متغيرًا من
  النمط~$\tau$، كان $\lambd[x][s]$ حدًا من النمط $\tau \to \sigma$.
\item إذا كان $s$ حدًا من النمط~$\sigma$ وكان $t$ حدًا من
  النمط~$\tau$، كان $\tuple{s, t}$ حدًا من النمط $\sigma \times
  \tau$.
\item إذا كان $s$ حدًا من النمط~$\sigma \times \tau$، كان $p_1(s)$
  حدًا من النمط~$\sigma$، و$p_2(s)$ حدًا من النمط~$\tau$.
\end{enumerate}
ويدل $\Obj{R}_{st}$ حدسيًا على الدالة المعرفة عوديًا بـ
\begin{align*}
\Obj{R}_{st}(0) & = s \\
\Obj{R}_{st}(x+1) & = t(x, R_{st}(x)),
\end{align*}
ويدل $\tuple{s, t}$ على الزوج الذي مركبته الأولى~$s$ ومركبته
الثانية~$t$، ويدل $p_1(s)$ و$p_2(s)$ على العنصرين الأول والثاني
(«الإسقاطين») من~$s$. وأخيرًا، يدل $\lambd[x][s]$ على الدالة~$f$
المعرفة بـ
\[
f(x) = s
\]
لكل~$x$ من النمط~$\tau$؛ ومن ثم يعطينا البند (6) صورة من الاستيعاب،
تمكننا من تعريف الدوال باستعمال الحدود. وتُبنى ال!!{formula}s من
عبارات ال!!{identity}~$\eq[s][t]$ بين حدود من النمط نفسه، والروابط
القضوية المعتادة، والتسوير على الأنماط العليا. ويمكن عندئذ اتخاذ
المعادلات الأساسية التي تحكم الحدود المعرفة أعلاه بديهيات للنسق، إلى
جانب قواعد المنطق المعتادة ذات المكممات وال!!{identity}.

وإذا زدنا نسق الأنماط المنتهية نمطًا~$\Omega$ لقيم الصدق، لزم أن ندرج
بديهيات تحكم استعماله أيضًا. بل يمكننا، إذا أحسنا التدبير، أن نستغني
تمامًا عن ال!!{formula}s المركبة ونستبدل بها حدودًا من النمط~$\Omega$!{}
ويمكن عندئذ تعديل نسق الإثبات وفقًا لذلك. والحاصل هو في جوهره
\emph{نظرية الأنماط البسيطة} التي قدمها ألونزو تشرش في ثلاثينيات القرن
العشرين.

وكما في منطق الرتبة الثانية، ثمة صور مختلفة لدلالة الأنماط العليا قد
نرغب في استعمالها. ففي الصورة التامة تتراوح متغيرات النمط $\sigma
\to \tau$ على مجموعة \emph{جميع} الدوال من أشياء النمط~$\sigma$ إلى
أشياء النمط~$\tau$. وكما قد تتوقع، هذه الدلالة أقوى من أن تقبل نسق
!!{derivation} فعالًا وتامًا. لكن يمكن النظر في دلالة أضعف، تتألف فيها
!!a{structure} من مجموعات عناصر~$T_\tau$ لكل نمط~$\tau$، مع العمليات
المناسبة للتطبيق والإسقاط وغيرهما. وإذا نُفذت التفاصيل على الوجه
الصحيح، أمكن الحصول على مبرهنات تمام لأنواع أنساق ال!!{derivation}
الموصوفة أعلاه.

يجذب منطق الأنماط العليا لأنه يقدم إطارًا نستطيع أن نضمّن فيه قدرًا
كبيرًا من الرياضيات على نحو طبيعي: فمن~$\Nat$ يمكن تعريف الأعداد
الحقيقية والدوال المتصلة، وهلم جرًا. وهو جذاب بوجه خاص في سياق المنطق
الحدسي، لأن للأنماط تأويلات «بنائية» واضحة. بل يمكن تطوير صور بنائية
من دلالة الأنماط العليا (قائمة على المنطق الحدسي لا الكلاسيكي) توضح
هذه التأويلات البنائية توضيحًا حسنًا، وتكون من وجوه كثيرة أشد إثارة من
نظائرها الكلاسيكية.

\end{document}
